% Propietario conservado: /Users/ruben/Documents/ChatGPT/jueces y controles/REVISION_ARGUMENTAL_APERTURAS_CIERRES_20260915/II/manuscript_es/sections/10d_portador_pentadico_rev06.tex
% RESULTADO_RECUPERADO. Integral activo, capítulos 63--65.
% B023_ch63_descenso_icosaedrico_portador_6ba97d5bb029_07_gravedad.tex:196--380.
% B024_ch64_entrelazador_tensorial_reduccion_6ba97d5bb029_07_gravedad.tex:1--122.
% B025_ch65_descomposicion_helicoidal_espin_dos_6ba97d5bb029_07_gravedad.tex:1--104.
% Adecuación editorial: jerarquía, referencias locales y distinción del proyector icosaédrico.
\subsection{Icosahedral incidence and five-component gravitational coupling}
\label{ii:sec:ii-pentadica}

The discrete structure of the continuum preserves an incidence reading
alongside the arithmetic reading. The five-component gravitational carrier is
constructed through that incidence and its tensor representation. The coupling
section \(G_a(\vartheta^\circ_{108})\), with
\(a\in\{\mathrm{pre},\mathrm{ret}\}\), is the one obtained in
section~\ref{ii:sec:gravity}; what follows makes explicit the space on
which it acts, its five components and its transverse projection.

Let \(X_{\rm TPK}^{\rm exc}\) be the subspace of the enriched state \(\TPK\)
that preserves, in addition to the arithmetic output, the orientation, route,
carry, and exceptional incidence. To make the geometric target space explicit,
let \(\varphi_g=(1+\sqrt5)/2\), \(\nu=\sqrt{1+\varphi_g^2}\), and define
\begin{equation}
\label{ii:ii:pent:eq:gravity-icosahedron-vertices}
\Omega_{12}^{\rm ico}
=\frac1\nu\bigl(
\{0\}\times\{\pm1\}\times\{\pm\varphi_g\}
\;\cup\;
\{\pm1\}\times\{\pm\varphi_g\}\times\{0\}
\;\cup\;
\{\pm\varphi_g\}\times\{0\}\times\{\pm1\}
\bigr)\subset S^2.
\end{equation}
Its twelve elements are the oriented vertex classes of an icosahedron.
The incidence is not left verbal: it is fixed by
\begin{equation}
\label{ii:ii:pent:eq:gravity-icosahedron-incidence}
\mathcal I_{12}^{\rm ico}
=\{(v,w)\in\Omega_{12}^{\rm ico}\times\Omega_{12}^{\rm ico}:
v\neq w,\ \langle v,w\rangle=1/\sqrt5\}.
\end{equation}
Each vertex has five neighbors in this relation. The reading bridge
is typed as a map
\begin{equation}
\label{ii:ii:pent:eq:gravity-exceptional-reader-map}
\pi_{\rm ico}:X_{\rm TPK}^{\rm exc}\longrightarrow\Omega_{12}^{\rm ico}.
\end{equation}
The incidence realization requires, in addition to the twelve elements of the
image, the following simultaneous properties of the map:
\begin{enumerate}[label=(\roman*)]
\item being surjective and constant exactly on the fibers that
discard the arithmetic chart but preserve the exceptional route;
\item intertwine the action of the oriented automorphisms of the antecedent
with a transitive action on \(\Omega_{12}^{\rm ico}\);
\item transporting the declared TPK incidence to
\(\mathcal I_{12}^{\rm ico}\), not
merely its cardinalities;
\item preserve the leaf involution and the nonadic return in the
stabilizer of the corresponding fiber.
\end{enumerate}

These conditions fix the order of the construction: first the map
\(\pi_{\rm ico}\), then the identification of the induced action and,
finally, its linear representation. Orientation, route and carry
belong to the enriched antecedent of that descent.


% BEGIN OWNER /Users/ruben/Documents/ChatGPT/jueces y controles/REVISION_ARGUMENTAL_APERTURAS_CIERRES_20260915/II/manuscript_es/sections/representacion_dodecafasica_rev08.tex
% FORMALIZACION_REUNIDA de la proyección del calendario y la representación
% de doce posiciones ya desarrolladas en la arquitectura TPK y el enlace Witt.
% Formalización explícita del cotejo de Gram y del transporte de marco.
\subsubsection{Calendar projection and representation of its twelve positions}
\label{ii:ii:dodeca:calendario}
The calendar of Section~\ref{ii:ii:tpk-arquitectura} provides an
observable coordinate of the complete state. In its notation, write
\(r_0(\theta)=r(\theta)-1\in\{0,\ldots,8\}\) and
\(b=\beta(\theta)\in\mathbb Z/12\mathbb Z\). The clock projection is
\[
 Q_{\rm clk}(x)=\bigl(\beta(\theta(x)),r_0(\theta(x))\bigr).
\]
For an admissible route of \(n\) updates, the induced transport is
\begin{equation}
 k(r_0,n)=\left\lfloor\frac{r_0+n}{9}\right\rfloor,
 \qquad
 T_n(b,r_0)=\bigl(b+k(r_0,n)\bmod12,\ (r_0+n)\bmod9\bigr).
 \label{ii:ii:dodeca:transporte-reloj}
\end{equation}
Here \(n\in\mathbb N_0\) counts effective updates; the integer
quotient \(k\) records the nonadic returns crossed.

\begin{proposition}[Compatibility of the calendar with route composition]
\label{ii:ii:dodeca:composicion}
The projection \(Q_{\rm clk}\), together with the length of each route, defines
a functor from the category of admissible TPK routes to the action category
of calendar~\eqref{ii:ii:dodeca:transporte-reloj}. The identities
\(T_{m+n}=T_mT_n\), \(T_0=I\),
\(T_9(b,r_0)=(b+1,r_0)\), and \(T_{108}(b,r_0)=(b,r_0)\) hold.
\end{proposition}
\begin{proof}
Putting \(r'_0=(r_0+n)\bmod9\), Euclidean division gives
\[
 k(r_0,n+m)=k(r_0,n)+k(r'_0,m).
\]
This identity simultaneously preserves residue and quotient when the
routes are concatenated. The update of \(\theta\) established in
Section~\ref{ii:ii:tpk-arquitectura} induces exactly
\eqref{ii:ii:dodeca:transporte-reloj}; the remaining identities follow
by substitution. The target category has pairs \((b,r_0)\) as
objects and admissible lengths as arrows, with addition as composition.
\end{proof}

The index \(p=b+1\in\{1,\ldots,12\}\), taking the representative
\(0\le b<12\), publishes the dodecaphase position. One turn of
\(108\) updates traverses its twelve positions. The proved periodicity
corresponds to this projection: route, accumulated carry and
memory remain in the source state and continue their update.

\paragraph{Incidence realization from the angular projector.}
The representatives \(r_pC_5\) and the representation \(\rho\) of
section~\ref{ii:sec:ii-enlace-angular-witt} fix the chart of the twelve
positions. The projector already defined by
\eqref{ii:eq:ii-witt-projector-character} produces
\begin{equation}
 \begin{gathered}
 \mathcal H_3=P_3^{\mathrm{ico}}\mathbb R^{12},\qquad v_p=2P_3^{\mathrm{ico}}e_p,
 \\[-2pt]
 (P_3^{\mathrm{ico}})_{pq}=\frac1{20}
 \sum_{\substack{g\in A_5\\gr_qC_5=r_pC_5}}\chi_3(g^{-1}).
 \end{gathered}
 \label{ii:ii:dodeca:vertices-proyector}
\end{equation}
The last sum contains five terms per entry and uses the
representatives and the two classes of order-five cycles fixed there.

\begin{proposition}[Icosahedral realization of the dodecaphase positions]
\label{ii:ii:dodeca:realizacion}
The vectors \(v_p\) are twelve distinct unit vectors in
\(\mathcal H_3\). Their incidence is
\(p\sim q\Longleftrightarrow\langle v_p,v_q\rangle=1/\sqrt5\).
The map \(p\mapsto v_p\) is \(A_5\)-equivariant, and its image
is isometric to \(\Omega_{12}^{\rm ico}\).
\end{proposition}
\begin{proof}
Transitivity and \(\operatorname{Tr}P_3^{\mathrm{ico}}=3\) give
\((P_3^{\mathrm{ico}})_{pp}=1/4\); by idempotence,
\(\langle v_p,v_q\rangle=4(P_3^{\mathrm{ico}})_{pq}\) and \(\|v_p\|=1\).
Evaluation of the five terms in
\eqref{ii:ii:dodeca:vertices-proyector} gives the antipodal pairs
\[
 (1,4),\ (2,3),\ (5,8),\ (6,7),\ (9,12),\ (10,11).
\]
For the representatives \((1,2,5,6,9,10)\), their Gram matrix is
\begin{equation}
 I_6+\frac D{\sqrt5},\qquad
 D=\begin{pmatrix}
 0&-1&-1&1&-1&1\\
 -1&0&1&1&-1&-1\\
 -1&1&0&1&1&1\\
 1&1&1&0&-1&1\\
 -1&-1&1&-1&0&1\\
 1&-1&1&1&1&0
 \end{pmatrix},\qquad D^2=5I_6.
 \label{ii:ii:dodeca:gram}
\end{equation}
The remaining scalar products are obtained by antipodality.
Thus, apart from a vertex and its opposite, the products are
\(\pm1/\sqrt5\), and each vertex has five neighbours. The following table
makes explicit the comparison with the signed six-position chart
\(\mathcal U=(\infty,0,1,2,3,4)\):
\begin{equation}
\begin{array}{c|cc@{\qquad}c|cc}
p&u(p)&\sigma(p)&p&u(p)&\sigma(p)\\ \hline
1&\infty&+1&7&3&-1\\
2&0&-1&8&1&+1\\
3&0&+1&9&2&-1\\
4&\infty&-1&10&4&+1\\
5&1&-1&11&4&-1\\
6&3&+1&12&2&+1
\end{array}
\label{ii:ii:dodeca:carta-firmada}
\end{equation}
Indeed, for \(C=\operatorname{bal}(A_W)\), with \(A_W\) given in
\eqref{ii:exc:aw}, the entries of \(D\) are
\(D_{ab}=\sigma(p_a)\sigma(p_b)C_{u(p_a),u(p_b)}\).
Consequently, the table and~\eqref{ii:ii:dodeca:gram} isometrically
identify \(v_p\) with
\(\sigma(p)\sqrt2 E_Ce_{u(p)}\), where
\(E_C=(I_6+C/\sqrt5)/2\). Both systems span three-dimensional
spaces and have identical Gram matrices, proving that the induced
linear map is well defined and isometric. The explicit
realization in~\eqref{ii:ii:ico:proyector-seis}, proved below,
identifies it with~\(\Omega_{12}^{\rm ico}\) and its incidence.
Finally, \(P_3^{\mathrm{ico}}\rho(g)=\rho(g)P_3^{\mathrm{ico}}\) implies
\(v_{g\cdot p}=\rho(g)v_p\). The bijection of the twelve positions gives
\(\operatorname{Stab}_{A_5}(v_p)=r_pC_5r_p^{-1}\).
\end{proof}

One thus obtains the representational composition
\begin{equation}
 x\longmapsto Q_{\rm clk}(x)\longmapsto
 p=\beta(\theta(x))+1\longmapsto v_p.
 \label{ii:ii:dodeca:composicion-representacional}
\end{equation}
The signed expression of this composition is
\(r_\pm^{\rm rep}(x)=(u(p),\sigma(p))\), determined by
table~\eqref{ii:ii:dodeca:carta-firmada}. The table is an explicit choice
of alignment between charts; the actions of \(A_5\) provide its
equivalent alignments. The complete register is preserved in the
domain of~\eqref{ii:ii:dodeca:composicion-representacional}.

\paragraph{Link with the five-component tensor carrier.}
In \(\mathcal H_3\), form
\(Q_p=v_pv_p^*-I_{\mathcal H_3}/3\). Let \(J=\mathbf1\mathbf1^*\),
and let \(\mathsf A\) be the matrix of the preceding antipodal permutation.
The Frobenius product and~\eqref{ii:ii:dodeca:gram} give
\begin{equation}
 \begin{aligned}
 \operatorname{Gram}(Q_p)
 &=(4P_3^{\mathrm{ico}})\circ(4P_3^{\mathrm{ico}})-\frac13J
 =\frac45(I+\mathsf A)-\frac2{15}J,
 \\
 P^{\rm ico}_5&=\frac58\operatorname{Gram}(Q_p)
 =\frac12(I+\mathsf A)-\frac1{12}J.
 \end{aligned}
 \label{ii:ii:dodeca:puente-tensorial}
\end{equation}
The symbol \(\circ\) denotes the entrywise product. The last
expression is the orthogonal projector onto coordinates even under
antipodality and orthogonal to \(\mathbf1\), of dimension five. Its character
is \((5,1,-1,0,0)\): on the six antipodal pairs, the fixed-point
counts are \((6,2,0,1,1)\), and the constant component is removed.
It therefore coincides with the central projector of
\eqref{ii:ii:pent:eq:gravity-p5}. This identity links the three-dimensional
construction with \(\Gamma_0\) and \(\Gamma_5\), defined in
\eqref{ii:ii:pent:eq:gravity-Gamma0} and~\eqref{ii:ii:pent:eq:gravity-gamma5},
through the explicit quadratic operation \(v_p\mapsto Q_p\).

\subsubsection{Transport of the representational frame}
\label{ii:ii:dodeca:marco}
Let \(s=(1\,2\,\cdots\,12)\), and let \(S_ce_p=e_{s(p)}\).
A nonadic calendar return takes \(p\) to \(s(p)\). The joint
transport of frame, representation, and projector is expressed by
\begin{equation}
 \begin{gathered}
 P_{3,k}^{\mathrm{ico}}=S_c^kP_3^{\mathrm{ico}}S_c^{-k},\qquad
 \rho^{(k)}(g)=S_c^k\rho(g)S_c^{-k},\\
 v^{(k)}_{s^k(p)}=2P_{3,k}^{\mathrm{ico}}e_{s^k(p)}=S_c^kv_p.
 \end{gathered}
 \label{ii:ii:dodeca:marco-movil}
\end{equation}
The last equality follows by substitution; orthogonality of
\(S_c\) preserves all inner products and transports
incidence. For a route of length \(n\), take
\(k=k(r_0,n)\). The carry identity in
Proposition~\ref{ii:ii:dodeca:composicion} proves compatibility of
these transports with concatenation. Although \(S_c^{12}=I\),
the twelve turns update the memory of the source route.

Equality~\eqref{ii:ii:dodeca:marco-movil} establishes covariance between
frames. The active action on a fixed incidence has additional
content: it requires identifying the transport of the complete state and its
action in that chart. In particular, \(s\) has order twelve, whereas
the orders of the elements of \(A_5\) are \(1,2,3,5\).
Likewise, the stabilizer of a vertex has order five, and that of its
antipodal pair has order ten; the involution of the signed chart
interchanges the two vertices. These distinctions preserve the meaning
of orientation, sheet and return in properties (i)--(iv): the
proved representational realization and the descent of the complete
exceptional route are related through those properties, with their
declared domains and fibres.

% END OWNER


% BEGIN OWNER /Users/ruben/Documents/ChatGPT/jueces y controles/REVISION_ARGUMENTAL_APERTURAS_CIERRES_20260915/II/manuscript_es/sections/incidencia_icosaedrica_rev07.tex
% FORMALIZACION_NUEVA de la incidencia firmada de la carta Paley preexistente.
% La construcción comienza en exc:aw; no identifica por decreto sus banderas
% con el cociente de todas las historias TPK.
\paragraph{Construction of the carrier from six-position incidence.}
\label{ii:ii:ico:construccion-firmada}
The matrix \(A_W\) of \eqref{ii:exc:aw}, which enters previously in the
angular transport of section~\ref{ii:sec:ii-enlace-angular-witt}, provides
an explicit antecedent of the icosahedral realization. Let
\(C=\operatorname{bal}(A_W)\), where
\(\operatorname{bal}(0)=0\), \(\operatorname{bal}(1)=1\) and
\(\operatorname{bal}(2)=-1\). The order of its coordinates is
\(\mathcal U=(\infty,0,1,2,3,4)\). The integer chart satisfies
\(C=C^{\mathsf T}\), \(C^2=5I_6\) and \(\operatorname{Tr}C=0\).
We define its signed covering by
\begin{equation}
 \mathcal D_C=\mathcal U\times\{+1,-1\},\qquad
 (u,\sigma)\sim_C(v,\tau)
 \ \Longleftrightarrow\ u\ne v\ \text{ and }\ \sigma\tau=C_{uv}.
 \label{ii:ii:ico:incidencia-firmada}
\end{equation}
Orientation doubles each coordinate; the matrix decides the incidences
between the two corresponding fibers. The involution of this covering
is \(\jmath_C(u,\sigma)=(u,-\sigma)\).

\begin{proposition}[Isometric realization of signed incidence]
\label{ii:ii:ico:realizacion-firmada}
The operator and vectors
\begin{equation}
 E_C=\frac12\left(I_6+\frac C{\sqrt5}\right),\qquad
 \mathcal H_C=\operatorname{im}E_C,\qquad
 v_{u,\sigma}=\sigma\sqrt2 E_Ce_u
 \label{ii:ii:ico:proyector-seis}
\end{equation}
realize \(\mathcal D_C\) as twelve distinct unit vectors in a
three-dimensional Euclidean space. There is an explicit isometry
\(T_C:\mathcal H_C\to\mathbb R^3\) for which
\(\iota_C(u,\sigma)=T_Cv_{u,\sigma}\) is a bijection onto
\(\Omega_{12}^{\rm ico}\), preserves incidence, and satisfies
\(\iota_C\jmath_C=-\iota_C\).
\end{proposition}
\begin{proof}
From \(C^2=5I_6\), one obtains \(E_C^2=E_C=E_C^*\); its trace is three.
The metric identities are
\[
 \langle v_{u,\sigma},v_{v,\tau}\rangle
 =\begin{cases}\sigma\tau,&u=v,\\
 \sigma\tau C_{uv}/\sqrt5,&u\ne v.
 \end{cases}
\]
Thus the twelve vectors are distinct, and their inner-product
\(1/\sqrt5\) incidence is exactly~\eqref{ii:ii:ico:incidencia-firmada}.
To make the isometry explicit, let the columns of \(B\), in order
\(\mathcal U\), be
\[
 B=\frac1\nu
 \begin{pmatrix}
 0&-1&-\varphi_g&0&\varphi_g&1\\
 -1&-\varphi_g&0&1&0&-\varphi_g\\
 -\varphi_g&0&-1&-\varphi_g&-1&0
 \end{pmatrix}.
\]
Substitution of \(\varphi_g^2=\varphi_g+1\) gives
\(B^*B=I_6+C/\sqrt5=2E_C\). Hence
\(T_C=B/\sqrt2\big|_{\mathcal H_C}\) is an isometry, and
\(T_Cv_{u,\sigma}=\sigma Be_u\). The six columns of \(B\) and their
opposites are precisely the twelve vertices of
\eqref{ii:ii:pent:eq:gravity-icosahedron-vertices}. The last identity
also proves compatibility of the involutions.
\end{proof}

The preceding dodecaphasic realization constructs the twelve vectors
directly from the projector \(P_3^{\mathrm{ico}}\) and the calendar positions.
The signed covering provides another coordinate expression of that
carrier. For the reader \(\pi_{\rm ico}\) in the descent theorem,
the relation between the two expressions is written
\begin{equation}
 X_{\rm TPK}^{\rm exc}\xrightarrow{\ r_\pm\ }
 \mathcal D_C\xrightarrow[\simeq]{\ \iota_C\ }
 \Omega_{12}^{\rm ico},\qquad
 \pi_{\rm ico}(x)=\sigma(x)Be_{u(x)},\quad
 r_\pm(x)=(u(x),\sigma(x)).
 \label{ii:ii:ico:factorizacion-lector}
\end{equation}
Since \(\iota_C\) is a bijection, this factorization determines
\(r_\pm=\iota_C^{-1}\pi_{\rm ico}\). Thus, \(r_\pm\) expresses
in the signed chart the reader being used: its introduction
does not add an independent hypothesis for constructing the dodecaphase
carrier. The preceding correspondence table gives that expression
for its twelve positions. Properties (i)--(iv) retain their
own scope when the realization is identified with the quotient
of the complete exceptional route. In particular, the joint transport
of frame, projector and incidence proves covariance between charts;
the preservation of a fixed incidence by an active transformation
is a distinct identity. Orientation and memory remain
in the source state during both readings.
A signed permutation that transports \(C\) jointly transports
\(E_C\), its vectors and its incidence; this covariance allows the charts
to be compared while preserving the orientation of each fibre.

% END OWNER


\begin{theorem}[Incidence Descent to the Icosahedral Quotient]
\label{ii:ii:pent:thm:gravity-exceptional-descent}
Suppose that \(\pi_{\rm ico}\) satisfies (i)--(iv) and that the oriented action
induced on \(\Omega_{12}^{\rm ico}\) has group
\(G_{\rm ico}\simeq A_5\). Then, for every class
\(v_0\in\Omega_{12}^{\rm ico}\),
\begin{equation}
\label{ii:ii:pent:eq:gravity-icosahedral-quotient}
H_{v_0}=\operatorname{Stab}_{G_{\rm ico}}(v_0)\simeq C_5,
\qquad
\Omega_{12}^{\rm ico}\simeq G_{\rm ico}/H_{v_0}\simeq A_5/C_5.
\end{equation}
The application
\begin{equation}
\label{ii:ii:pent:eq:gravity-coset-map}
\kappa_{v_0}:G_{\rm ico}/H_{v_0}\longrightarrow\Omega_{12}^{\rm ico},
\qquad gH_{v_0}\longmapsto g\cdot v_0,
\end{equation}
is an equivariant bijection and preserves incidence if, in the quotient,
one declares
\begin{equation}
\label{ii:ii:pent:eq:gravity-coset-incidence}
gH_{v_0}\sim g'H_{v_0}
\quad\Longleftrightarrow\quad
(g\cdot v_0,g'\cdot v_0)\in\mathcal I_{12}^{\rm ico}.
\end{equation}
Therefore, the composite \(\kappa_{v_0}^{-1}\pi_{\rm ico}\) is the typed map from the TPK
antecedent to the quotient \(A_5/C_5\).
\end{theorem}

\begin{proof}
The list of \eqref{ii:ii:pent:eq:gravity-icosahedron-vertices} has twelve elements. By transitivity, \(|G_{\rm ico}\cdot v_0|=12\), and the orbit--stabilizer theorem gives \(|H_{v_0}|=60/12=5\). Every group of prime order five is cyclic; hence \(H_{v_0}\simeq C_5\). The formula for \(\kappa_{v_0}\) does not depend on the representative: if \(gH_{v_0}=g'H_{v_0}\), then \(g^{-1}g'\in H_{v_0}\) and \(g'v_0=gv_0\). Surjectivity follows from transitivity, and equality of the cardinalities gives injectivity. Equivariance is immediate: \(\kappa_{v_0}(agH_{v_0})=a\kappa_{v_0}(gH_{v_0})\). Finally, \eqref{ii:ii:pent:eq:gravity-coset-incidence} is precisely the transport of structure through that bijection; hypothesis (iii) makes this transport commute with \(\pi_{\rm ico}\).
\end{proof}

The theorem establishes descent once the four
properties of \(\pi_{\rm ico}\) have been verified. Its domain is the enriched state
of the exceptional reading, which preserves the incidence transported by
APP--TRIT--TPK. In this construction, the explicit antecedent is the
map~\eqref{ii:ii:pent:eq:gravity-exceptional-reader-map} with its four
properties; the passage from \(A_5/C_5\) to the transverse traceless
reduction is proved below with that antecedent and the induced action
of \(A_5\).

The chain of provenance is expressed through all its morphisms:
\begin{equation}
\label{ii:ii:pent:eq:gravity-exceptional-provenance}
\APP\longrightarrow\mathrm{TRIT}\longrightarrow\TPK
\longrightarrow X_{\rm TPK}^{\rm exc}
\xrightarrow{\ \pi_{\rm ico}\ }\Omega_{12}^{\rm ico}
\xleftarrow[\simeq]{\ \kappa_{v_0}\ }A_5/C_5.
\end{equation}
The composition~\eqref{ii:ii:pent:eq:gravity-exceptional-provenance} retains
as its antecedent the existence of \(\pi_{\rm ico}\) with properties
(i)--(iv). On that domain, descent to the quotient is
theorem~\ref{ii:ii:pent:thm:gravity-exceptional-descent}; from it are constructed
the projector \(P^{\mathrm{ico}}_5\), the intertwiner \(\Gamma_5\) and the
projection \(\Lambda(n)\). The realization of the carrier as a gravitational
field additionally uses the dynamics and physical constraints
specified at the end of this development.

Let \(\mathcal G=A_5\), let \(H\simeq C_5\) be the stabilizer of an oriented vertex
of the icosahedron, and consider
\begin{equation}
\label{ii:ii:pent:eq:gravity-v12}
V_{12}=\R[\mathcal G/H],
\qquad |\mathcal G/H|=\frac{60}{5}=12.
\end{equation}

% BEGIN OWNER /Users/ruben/Documents/ChatGPT/jueces y controles/REVISION_ARGUMENTAL_APERTURAS_CIERRES_20260915/II/manuscript_es/sections/caracteres_pentadicos_rev07.tex
% FORMALIZACION_NUEVA: prueba explícita de la descomposición y la
% irreducibilidad ya enunciadas en el desarrollo pentacomponente de REV06.
% CERTIFICADO_NUEVO: verificar_caracteres_pentadicos_rev07.py.
% Inserción prevista tras la definición de V12, antes de su descomposición.
\paragraph{Characters of icosahedral incidence.}
The following calculation is performed on the icosahedral action obtained in
Theorem~\ref{ii:ii:pent:thm:gravity-exceptional-descent}, with its incidence
and orientation hypotheses. It determines the decomposition of the
vertex representation and the irreducibility of the tensor space
through their traces, preserving the antecedent
\(\pi_{\mathrm{ico}}\) declared in that theorem.

\begin{proposition}[Vertex decomposition and tensor character]
\label{ii:prop:ii-pent-caracteres}
Let \(\tau=(1+\sqrt5)/2\) and
\(\bar\tau=(1-\sqrt5)/2\). In the class order
\(1,2A,3A,5A,5B\), the characters of the rotational representation
\(V_3\), its algebraic conjugate \(V_{3'}\), the twelve-vertex
representation \(V_{12}\), and
\(\operatorname{STF}_3=\operatorname{Sym}^2_0(V_3)\) are
\begin{equation}
\label{ii:eq:ii-pent-tabla-caracteres}
\begin{array}{c|rrrrr}
\text{class} &1&2A&3A&5A&5B\\ \hline
\text{cardinal}&1&15&20&12&12\\
\text{class of }g^2&1&1&3A&5B&5A\\ \hline
\chi_3&3&-1&0&\tau&\bar\tau\\
\chi_{3'}&3&-1&0&\bar\tau&\tau\\
\chi_{12}&12&0&0&2&2\\
\chi_{\mathrm{STF}}&5&1&-1&0&0
\end{array}
\end{equation}
The representations \(V_3,V_{3'}\) and
\(\operatorname{STF}_3\) are absolutely irreducible and distinct.
Moreover,
\begin{equation}
\label{ii:eq:ii-pent-descomposicion-demostrada}
 V_{12}\simeq\mathbf1\oplus V_3\oplus V_{3'}
                  \oplus\operatorname{STF}_3,
 \qquad
 \langle\chi_{12},\chi_{\mathrm{STF}}\rangle=1.
\end{equation}
In particular, the five-dimensional summand occurs exactly once.
\end{proposition}

\begin{proof}
In \(A_5\), there are \(5\cdot3=15\) double transpositions and
\(\binom53\cdot2=20\) three-cycles.
The \(4!=24\) five-cycles split into two classes of
twelve: their centralizer is the cyclic subgroup of order five and is
contained in \(A_5\). Identifying the positions of a cycle
with \(\mathbb Z/5\mathbb Z\), the permutation inducing
\(r\mapsto2r\) is a four-cycle on the nonzero
elements, with negative sign. Squaring therefore exchanges the two
classes. Inversion \(r\mapsto-r\) is a product of two
transpositions and preserves each class. This determines the first two
rows of~\eqref{ii:eq:ii-pent-tabla-caracteres}.

The representation \(V_3\) is the action by rotations on the
vertex coordinates. A rotation through angle \(\theta\)
has trace \(1+2\cos\theta\). The classes of orders two and three
give \(-1\) and \(0\), respectively. The name \(5A\) denotes the
class of angles \(\pm2\pi/5\) and \(5B\) that of
\(\pm4\pi/5\); their traces are \(\tau\) and \(\bar\tau\).
These expressions follow from
\(u^2+u-1=0\), for \(u=2\cos(2\pi/5)\), and from the double-angle
formula. The equation for \(u\) follows by dividing
\(1+z+z^2+z^3+z^4=0\) by \(z^2\), with
\(z=\exp(2\pi i/5)\). We have thus calculated the character of a
representation already defined geometrically.

To construct the other representation, use the vertices without the
common normalization \(\nu^{-1}\). Their coordinates belong to
\(K=\mathbb Q(\sqrt5)\). If \(B\) is the matrix of three linearly
independent vertices and \(B_g\) contains their images, the
rotation matrix is \(R_g=B_gB^{-1}\in M_3(K)\).
The conjugation \(\sigma:\sqrt5\mapsto-\sqrt5\) then defines
\(R'_g=\sigma(R_g)\), satisfying
\[
 R'_{gh}=R'_gR'_h,\qquad
 (R'_g)^TR'_g=I,\qquad \det R'_g=1.
\]
Thus \(V_{3'}\) is an effectively constructed real
representation. Its character is \(\sigma(\chi_3)\), exchanging
\(\tau\) and \(\bar\tau\).

The permutation character counts fixed vertices. The identity fixes
twelve; nontrivial rotations of order five fix the two endpoints
of their axis. Those of order two or three fix no vertices, since the
stabilizer of a vertex is \(C_5\). Hence
\(\chi_{12}=(12,0,0,2,2)\).

For the tensor space, if \(\lambda_1,\lambda_2,\lambda_3\)
are the eigenvalues of \(R_g\), the sum
\(\sum_{i\leq j}\lambda_i\lambda_j\) is
\(\bigl((\sum_i\lambda_i)^2+\sum_i\lambda_i^2\bigr)/2\).
The metric form generates an invariant line in
\(\operatorname{Sym}^2(V_3)\), whose trace-free complement has
character
\begin{equation}
\label{ii:eq:ii-pent-caracter-stf}
 \chi_{\mathrm{STF}}(g)
 =\frac{\chi_3(g)^2+\chi_3(g^2)}2-1.
\end{equation}
Using the squaring map in the table and
\(\tau+\bar\tau=1\), \(\tau\bar\tau=-1\), gives
\(\chi_{\mathrm{STF}}=(5,1,-1,0,0)\).

The real-character inner product is
\(\langle\chi,\psi\rangle
=\frac1{60}\sum_C|C|\chi(C)\psi(C)\).
In particular,
\begin{align*}
 \langle\chi_3,\chi_3\rangle
 &=\langle\chi_{3'},\chi_{3'}\rangle
 =\frac{9+15+12(\tau^2+\bar\tau^2)}{60}=1,\\
 \langle\chi_3,\chi_{3'}\rangle
 &=\frac{9+15+24\tau\bar\tau}{60}=0,\\
 \langle\chi_{\mathrm{STF}},\chi_{\mathrm{STF}}\rangle
 &=\frac{25+15+20}{60}=1.
\end{align*}
Norm one proves irreducibility of the complexifications;
consequently, it also proves absolute irreducibility of the constructed
real representations. Their dimensions and orthogonality distinguish the
three summands. Finally, the equality of class functions
\[
 \chi_{12}=1+\chi_3+\chi_{3'}+\chi_{\mathrm{STF}}
\]
is verified in all five columns. Semisimplicity of finite-group
representations in characteristic zero gives
decomposition~\eqref{ii:eq:ii-pent-descomposicion-demostrada}. The
multiplicity of the tensor summand can be checked directly:
\[
 \langle\chi_{12},\chi_{\mathrm{STF}}\rangle
 =\frac{12\cdot5}{60}=1.
\]
\end{proof}

Thus the central idempotent associated with
\(\chi_{\mathrm{STF}}=\chi_5\) has rank
\(5\cdot1=5\) in \(V_{12}\). The irreducibility entering
normalization of the tensor intertwiner and the multiplicity
entering gravitational coupling are established by
this same calculation.

% END OWNER


The permutation representation \(\rho_{12}\) decomposes as
\begin{equation}
\label{ii:ii:pent:eq:gravity-v12-decomposition}
V_{12}\simeq\mathbf1\oplus V_3\oplus V_{3'}\oplus V_5.
\end{equation}
The central idempotent of the character summand
\(\chi_5=(5,1,-1,0,0)\) is
\begin{equation}
\label{ii:ii:pent:eq:gravity-p5}
\boxed{
P^{\mathrm{ico}}_5=\frac5{60}\sum_{g\in A_5}\chi_5(g^{-1})\rho_{12}(g)
=\frac1{12}\left(
5I+\sum_{g\in2A}\rho_{12}(g)
-\sum_{g\in3A}\rho_{12}(g)
\right).}
\end{equation}

\begin{theorem}[Rank-five gravitational projector and coupler]
\label{ii:ii:pent:thm:gravity-p5-coupling}
The operator \(P^{\mathrm{ico}}_5\) satisfies
\begin{equation}
\label{ii:ii:pent:eq:gravity-p5-properties}
(P^{\mathrm{ico}}_5)^2=P^{\mathrm{ico}}_5=(P^{\mathrm{ico}}_5)^*,
\qquad \Tr P^{\mathrm{ico}}_5=\operatorname{rank} P^{\mathrm{ico}}_5=5.
\end{equation}
For
\begin{equation}
\label{ii:ii:pent:eq:gravity-G5}
\mathbb G_{5,a}:=G_a(\vartheta^\circ_{108})P^{\mathrm{ico}}_5
\end{equation}
we have
\begin{equation}
\label{ii:ii:pent:eq:gravity-G5-spectrum}
\Spec(\mathbb G_{5,a})
=\{G_a(\vartheta^\circ_{108})\text{ with multiplicity }5,
\;0\text{ with multiplicity }7\}.
\end{equation}
\end{theorem}

\begin{proof}
Character orthogonality converts \eqref{ii:ii:pent:eq:gravity-p5} into the
central idempotent that acts as the identity on the unique copy of
\(V_5\) and as zero on the other three summands of
\eqref{ii:ii:pent:eq:gravity-v12-decomposition}. Reality of the character and
orthogonality of \(\rho_{12}\) give self-adjointness. Its trace and rank
are five. Multiplying a rank-five projector by the positive scalar
\(G_a\) immediately produces the indicated spectrum.
\end{proof}

The gravitational section \(G_a\) is common to the five components of the
selected tensor summand. Its spectral multiplicity expresses the
rank of the carrier; the other seven modes of the permutation module
belong to the kernel of \(\mathbb G_{5,a}\). The symbol
\(P^{\mathrm{ico}}_5\) identifies here this rank-five central projector
on \(V_{12}\); the average of the five positions of an
arithmetic pentafibre has another domain and rank one.

\subsection{Intertwiner with the space of symmetric traceless tensors}

Let \(\mathcal V=\Omega_{12}^{\rm ico}\subset S^2\) be the already constructed set of
twelve unit vertices. For each \(v\in\mathcal V\), define
\begin{equation}
\label{ii:ii:pent:eq:gravity-Qv}
Q_v=vv^T-\frac13I_3\in
\operatorname{STF}_3:=\operatorname{Sym}^2_0(\R^3).
\end{equation}
The frame map
\begin{equation}
\label{ii:ii:pent:eq:gravity-Gamma0}
\Gamma_0\!\left(\sum_{v\in\mathcal V}x_ve_v\right)
=\sum_{v\in\mathcal V}x_vQ_v
\end{equation}
is \(A_5\)-equivariant and satisfies
\begin{equation}
\label{ii:ii:pent:eq:gravity-tight-frame}
\Gamma_0\Gamma_0^*=\frac85I_{\operatorname{STF}_3},
\qquad
\Gamma_0=\Gamma_0P^{\mathrm{ico}}_5.
\end{equation}

\begin{theorem}[Normalized tensor entangler]
\label{ii:ii:pent:thm:gravity-gamma5}
The map
\begin{equation}
\label{ii:ii:pent:eq:gravity-gamma5}
\boxed{
\Gamma_5:=\sqrt{\frac58}\,\Gamma_0|_{V_5}:
V_5\xrightarrow{\ \simeq\ }\operatorname{STF}_3}
\end{equation}
is a \(A_5\)-equivariant isometry. Its adjoint is
\begin{equation}
\label{ii:ii:pent:eq:gravity-gamma5-adjoint}
\Gamma_5^*Q=\sqrt{\frac58}
\sum_{v\in\mathcal V}(v^TQv)e_v.
\end{equation}
\end{theorem}

\begin{proof}
Equivariance follows from \(Q_{gv}=gQ_vg^T\). The operator
\(\Gamma_0\Gamma_0^*\) commutes with the irreducible representation of
\(A_5\) on \(\operatorname{STF}_3\), and is therefore scalar. Its trace
is
\[
\sum_{v\in\mathcal V}\|Q_v\|_F^2
=12\cdot\frac23=8;
\]
dividing by dimension five yields \(8/5\). Hence,
\(\Gamma_5\Gamma_5^*=I\). Since domain and codomain have dimension
five, also \(\Gamma_5^*\Gamma_5=I\). The formula for the adjoint
follows from \(\langle Q,Q_v\rangle_F=v^TQv\).
\end{proof}

Normalization through the icosahedral frame constructs the isometry that
transports the finite summand to the real spin-two tensor space.

\subsection{Transverse traceless projection}

For \(n\in S^2\), let \(\Pi(n)=I-nn^T\). We define
\begin{equation}
\label{ii:ii:pent:eq:gravity-lambda-tt}
\boxed{
\Lambda(n)Q
=\Pi Q\Pi-\frac12\Tr(\Pi Q\Pi)\Pi.}
\end{equation}

\begin{theorem}[Reduction from five components to two]
\label{ii:ii:pent:thm:gravity-tt}
The restriction of \(\Lambda(n)\) to \(\operatorname{STF}_3\) is the
orthogonal projector onto
\begin{equation}
\label{ii:ii:pent:eq:gravity-htt}
\mathcal H_{\rm TT}(n)
=\{Q\in\operatorname{STF}_3:Qn=0\}.
\end{equation}
In particular,
\begin{equation}
\label{ii:ii:pent:eq:gravity-lambda-properties}
\Lambda(n)^2=\Lambda(n)=\Lambda(n)^*,
\qquad \operatorname{rank}\Lambda(n)=2.
\end{equation}
For \(n=e_3\), the image is generated by
\begin{equation}
\label{ii:ii:pent:eq:gravity-plus-cross}
E_+=\frac1{\sqrt2}\operatorname{diag}(1,-1,0),
\qquad
E_\times=\frac1{\sqrt2}
\begin{pmatrix}0&1&0\\1&0&0\\0&0&0\end{pmatrix}.
\end{equation}
\end{theorem}

\begin{proof}
Since \(\Pi n=0\), the output is transverse. Its trace is
\[
\Tr(\Pi Q\Pi)-\frac12\Tr(\Pi)\Tr(\Pi Q\Pi)=0,
\]
since \(\Tr\Pi=2\). If \(Qn=0\) and \(\Tr Q=0\), then
\(\Pi Q\Pi=Q\) and \(\Tr(\Pi Q\Pi)=0\), so that
\(\Lambda(n)Q=Q\). The formula is self-adjoint with respect to the
Frobenius product. In the frame \(n=e_3\), its matrix in an adapted STF basis is
\(\operatorname{diag}(1,1,0,0,0)\), which proves the rank.
\end{proof}

The complete chain thus remains,
\begin{equation}
\label{ii:ii:pent:eq:gravity-12-5-5-2}
\boxed{
V_{12}\xrightarrow{P^{\mathrm{ico}}_5}V_5
\xrightarrow{\Gamma_5}\operatorname{STF}_3
\xrightarrow{\Lambda(n)}\mathcal H_{\rm TT}(n),
\qquad 12\longrightarrow5\longrightarrow5\longrightarrow2.}
\end{equation}
The reduced coupling
\begin{equation}
\label{ii:ii:pent:eq:gravity-effective-map}
\mathbb G_{{\rm TT},a}(n)
=G_a(\vartheta^\circ_{108})\Lambda(n)\Gamma_5P^{\mathrm{ico}}_5
\end{equation}
has rank two.

\subsection{Decomposition into \(5\) polarization modes}

Dimension five acquires tensorial content upon choosing a unit direction
of propagation \(n\) and an oriented orthonormal frame
\((e_1,e_2,n)\). Define the five real tensors
\begin{align}
E_{+}&=\frac1{\sqrt2}(e_1e_1^T-e_2e_2^T),
&E_{\times}&=\frac1{\sqrt2}(e_1e_2^T+e_2e_1^T),
\label{ii:ii:pent:eq:gravity-helicity2-real}\\
E_{1}&=\frac1{\sqrt2}(e_1n^T+ne_1^T),
&E_{2}&=\frac1{\sqrt2}(e_2n^T+ne_2^T),
\label{ii:ii:pent:eq:gravity-helicity1-real}\\
E_{0}&=\frac1{\sqrt6}(2nn^T-e_1e_1^T-e_2e_2^T).
\label{ii:ii:pent:eq:gravity-helicity0-real}
\end{align}
All are symmetric and traceless. The first two are transverse; the
next two mix the transverse plane with the direction of
propagation, and the last is the traceless longitudinal scalar mode.

\begin{theorem}[Complete helical decomposition of the carrier]
\label{ii:ii:pent:thm:gravity-five-helicities}
The tensors of
\eqref{ii:ii:pent:eq:gravity-helicity2-real}, \eqref{ii:ii:pent:eq:gravity-helicity1-real}, \eqref{ii:ii:pent:eq:gravity-helicity0-real}
form an orthonormal basis of \(\operatorname{STF}_3\). In the
complexification, the combinations
\begin{equation}
\label{ii:ii:pent:eq:gravity-complex-helicities}
E_{\pm2}=\frac1{\sqrt2}(E_+\mp i E_\times),
\qquad
E_{\pm1}=\frac1{\sqrt2}(E_1\mp i E_2),
\qquad E_0=E_0
\end{equation}
are weight-\(\pm2,\pm1,0\) vectors, respectively, under rotations
about \(n\). Moreover,
\begin{equation}
\label{ii:ii:pent:eq:gravity-tt-on-five}
\Lambda(n)E_{\pm2}=E_{\pm2},
\qquad
\Lambda(n)E_{\pm1}=0,
\qquad
\Lambda(n)E_0=0.
\end{equation}
\end{theorem}

\begin{proof}
In the frame \((e_1,e_2,n)\), each tensor has Frobenius norm one.
Their off-diagonal entry patterns are disjoint, and the three diagonals
of \(E_+,E_0\) have zero inner product; therefore, the five tensors
are orthonormal. Since \(\dim\operatorname{STF}_3=6-1=5\), they form a basis.

A rotation through angle \(\psi\) about \(n\) sends
\((e_1,e_2)\) to
\((\cos\psi\,e_1+\sin\psi\,e_2,
-\sin\psi\,e_1+\cos\psi\,e_2)\). Substituting into the formulas gives
a rotation through angle \(2\psi\) on
\(\operatorname{span}\{E_+,E_\times\}\), one through angle \(\psi\) on
\(\operatorname{span}\{E_1,E_2\}\), and \(E_0\) fixed. Hence the weights of
\eqref{ii:ii:pent:eq:gravity-complex-helicities}. Finally,
\(\Pi E_{+}\Pi=E_+\) and \(\Pi E_\times\Pi=E_\times\), whereas
\(\Pi E_1\Pi=\Pi E_2\Pi=0\). For \(E_0\),
\(\Pi E_0\Pi=-(e_1e_1^T+e_2e_2^T)/\sqrt6\), which is pure trace in
the plane, and subtracting \eqref{ii:ii:pent:eq:gravity-lambda-tt} annihilates it.
\end{proof}

The theorem specifies the meaning of ``five polarizations'' without assigning all of them the same dynamic status. Before field equations are imposed, the spin-two carrier has five irreducible coordinates. In a massive Fierz--Pauli realization, the constraints leave precisely the five weights of \eqref{ii:ii:pent:eq:gravity-complex-helicities}; in a massless realization with gauge invariance, the sectors \(\pm1\) and \(0\) are eliminated, leaving \(\pm2\). Here HMT supplies the carrier, the basis, and the projector; the choice of massive or massless dynamics belongs to the physical interface. Thus five does not mean ``five observed waves,'' but rather five tensorial degrees available before reduction.

The orientation also remains visible: replacing
\((e_1,e_2,n)\) with \((e_1,-e_2,-n)\) preserves \(E_+,E_0\) and exchanges
the phases of the helical pairs. That information is not in the
scalar \(G_a\); it resides in the carrier on which
\(\mathbb G_{5,a}\) acts.

\subsection{Physical interpretation of the representational chain}

Dimension five appears in distinct objects, and only
\eqref{ii:ii:pent:eq:gravity-12-5-5-2} authorizes transport between them:

\begin{enumerate}
\item \(V_5\) is the irreducible summand of the representation on the
twelve vertices. This is a finite result in representation theory.
\item \(\operatorname{STF}_3\) is the real five-dimensional carrier of
the \(m=-2,-1,0,1,2\) weights of spin two.
\item A massive Fierz--Pauli theory can realize these five weights
as five physical states, after fixing the action and constraints.
\item Massless general relativity retains two radiative helicities,
\(\pm2\), represented by \(E_+\) and \(E_\times\).
\item The \(E(2)\) classes of metric waves, such as \(III_5\), classify
possible amplitudes under other hypotheses; they are not automatically
\(V_5\), Fierz--Pauli, or general relativity.
\end{enumerate}

The structural content is thus determined: under the incidence-descent
conditions, HMT constructs a five-component tensor carrier
before imposing constraints. The massless realization preserves
the two radiative helicities through the rank-two TT projection.
Orientation information is preserved in the tensors and in their
transport; the scalar value \(G_a\) determines the common coupling.
